\documentclass[11pt,a4paper]{article}

\usepackage[T1]{fontenc}
\usepackage[utf8]{inputenc}
\usepackage[ngerman]{babel}
\usepackage{lmodern}
\usepackage{microtype}
\usepackage{amsmath,amssymb,amsthm,mathtools}
\usepackage{booktabs}
\usepackage{array}
\usepackage{enumitem}
\usepackage{tikz}
\usepackage{hyperref}
\usepackage{xcolor}

\hypersetup{
  colorlinks=true,
  linkcolor=blue!45!black,
  citecolor=blue!45!black,
  urlcolor=blue!45!black,
  pdftitle={Eine Kreis-Sehnen-Kodierung von Primzahlvierlingen},
  pdfsubject={Primzahlvierlinge, Restklassen, Kleinsche Vierergruppe, Kreis-Sehnen-Decoder},
  pdfkeywords={Primzahlvierlinge, prime quadruplets, Dirichlet-Charaktere, Kleinsche Vierergruppe, Restklassen, Lean}
}

\setlength{\parindent}{0pt}
\setlength{\parskip}{0.55em}
\setlist[itemize]{topsep=0.3em,itemsep=0.15em}
\setlist[enumerate]{topsep=0.3em,itemsep=0.15em}

\newtheorem{satz}{Satz}[section]
\newtheorem{lemma}[satz]{Lemma}
\newtheorem{korollar}[satz]{Korollar}
\newtheorem{proposition}[satz]{Proposition}
\theoremstyle{definition}
\newtheorem{definition}[satz]{Definition}
\newtheorem{bemerkung}[satz]{Bemerkung}
\theoremstyle{remark}
\newtheorem{beispiel}[satz]{Beispiel}

\newcommand{\Z}{\mathbb Z}
\newcommand{\N}{\mathbb N}
\newcommand{\F}{\mathbb F}
\newcommand{\eps}{\varepsilon}
\newcommand{\modulo}[1]{\; (\mathrm{mod}\,#1)}
\newcommand{\ABCE}{\mathrm{ABCE}}
\newcommand{\CEAB}{\mathrm{CEAB}}
\newcommand{\Vfour}{V_4}
\DeclareMathOperator{\Prime}{Prime}

\title{\textbf{Eine Kreis-Sehnen-Kodierung von Primzahlvierlingen}\\
\large Lokale Restklassen, Kleinsche Vierergruppe, dyadischer Decoder und eine Lean-formalisierbare Zerlegung}
\author{}
\date{7. Oktober 2026}

\begin{document}
\maketitle

\begin{abstract}
Wir untersuchen Primzahlvierlinge der Form
\[
Q(p)=(p,p+2,p+6,p+8)
\]
über ihre Zentralzahl $M=p+4$. Für alle nicht-exzeptionellen Vierlinge $p>5$ wird elementar gezeigt, dass
\[
M\equiv 15\pmod{30},\qquad M\pmod 7\in\{-1,0,+1\},
\]
und damit genau drei Restklassen $M\equiv15,105,195\pmod{210}$ auftreten. Die vier Primstellen des Vierlings liegen modulo $12$ in den vier Einheitenklassen $\{1,5,7,11\}$, die als Kleinsche Vierergruppe $(\Z/12\Z)^\times\simeq C_2\times C_2$ aufgefasst werden. Mit zwei reellen Dirichlet-Charakteren modulo $4$ und $3$ zerfällt diese Viererstruktur in einen starren Positionskanal und ein einziges mod-$4$-Orientierungsbit; dieses unterscheidet die beiden Orientierungen $\ABCE$ und $\CEAB$.

Als geometrisch-kominatorische Darstellung führen wir einen regulären Kreis mit $2M+1$ beschrifteten Punkten $-M,\ldots,0,\ldots,M$ ein und summieren gerichtete Differenzen über Parallelklassen von Sehnen. Für die resultierende Observable $D_M(c)$ wird die geschlossene Formel
\[
D_M(c)=M^2-(c-1)M-\Big\lfloor\frac c2\Big\rfloor
\]
hergeleitet. Insbesondere gilt die exakte Decoderidentität
\[
M=D_M(2)-D_M(3).
\]
Die drei mod-$210$-Zellen werden durch das mod-$4$-Orientierungsbit in sechs mod-$420$-Zellen aufgespalten. Ein allgemeiner dyadischer Lift zeigt, wie gerade und ungerade Sehnenrichtungen sukzessive zusätzliche $2$-adische Restklassenbits unterscheiden. Die Konstruktion liefert keine neue Primzahlexistenz- oder Dichtenaussage; sie ist eine invertierbare geometrisch-algebraische Kodierung bereits vorhandener lokaler Restklasseninformation. Abschließend werden eine Charakterprojektion, numerische Beispiele bis $10^7$ und ein Lean-4-Formalisierungsschema angegeben.
\end{abstract}

\medskip
\noindent\textbf{Schlüsselwörter:} Primzahlvierlinge; Primzahlkonstellationen; Restklassen; Dirichlet-Charaktere; Kleinsche Vierergruppe; Siebbedingungen; Kreis-Sehnen-Kodierung; Lean 4.\\
\textbf{MSC 2020:} 11A07, 11N05, 11N13, 20K01, 68V20.

\tableofcontents

\section{Einleitung und Geltungsbereich}

Primzahlvierlinge
\[
(p,p+2,p+6,p+8)
\]
gehören zu den dichtesten zulässigen Primzahlkonstellationen. Die lokale Zulässigkeit solcher Muster ist ein klassischer Ausgangspunkt der Theorie der Primzahl-$k$-Tupel: Eine endliche Verschiebungsmenge $\mathcal H$ kann nur dann unendlich oft vollständig prim auftreten, wenn sie modulo keiner Primzahl sämtliche Restklassen überdeckt. Die Hardy--Littlewood-Heuristik ordnet einer solchen Konstellation eine singuläre Reihe zu und sagt eine asymptotische Häufigkeit der Größenordnung $x/(\log x)^4$ voraus \cite{HardyLittlewood1923,IwaniecKowalski2004}. Für die spezielle Konstellation $\{0,2,6,8\}$ ist die Unendlichkeit der Primzahlvierlinge weiterhin unbewiesen.

Das Ziel dieser Arbeit ist bescheidener und zugleich strukturell: Wir isolieren die lokale Restklasseninformation eines Primzahlvierlings über seine Mitte
\[
M=p+4,
\]
zerlegen die vier Primstellen modulo $12$ in die Kleinsche Vierergruppe und geben eine Kreis-Sehnen-Kodierung an, aus der $M$ durch die Differenz zweier benachbarter Parallelklassensummen exakt zurückgewonnen wird.

Die Bezeichnungen $E,A,B,C$ sind \emph{keine Standardnotation der Zahlentheorie}; sie werden in diesem Manuskript als Kürzel für die vier Einheitenklassen modulo $12$ eingeführt. Ebenso ist der Ausdruck \emph{Kreis-Sehnen-Decoder} eine Bezeichnung für die hier definierte kombinatorische Kodierung.

Die zentrale methodische Grenze ist wichtig: Der Decoder ist invertierbar. Er kann daher keine zusätzliche arithmetische Information erzeugen, die nicht bereits in $M$ enthalten ist. Sein Nutzen liegt in der kanonischen Trennung verschiedener lokaler Kanäle und in einer geometrisch anschaulichen Darstellung der mod-$210$-, mod-$420$- und dyadischen Restklassenhierarchie.

\section{Historische und mathematische Einordnung}

\subsection{Von Dirichlet zu Charakteren}

Dirichlets Beweis von 1837 über unendlich viele Primzahlen in jeder reduzierten arithmetischen Progression markiert einen Ausgangspunkt der analytischen Zahlentheorie \cite{Dirichlet1837}. In diesem Zusammenhang traten jene multiplikativen Restklassenfunktionen auf, die heute als Dirichlet-Charaktere bezeichnet werden. Für die vorliegende Konstruktion benötigen wir keine analytische Theorie der $L$-Funktionen; wir verwenden lediglich zwei elementare reelle Charaktere auf den Einheiten modulo $12$.

Der historische Bezug ist dennoch sachlich: Die Zerlegung von Primzahlen nach Restklassen wird hier in derselben Sprache von Charakteren organisiert, die Dirichlet für die Untersuchung von Primzahlen in Progressionen eingeführt hat.

\subsection{Die Kleinsche Vierergruppe}

Die Gruppe mit vier Elementen, in der jedes von der Identität verschiedene Element Ordnung zwei besitzt, ist heute als Kleinsche Vierergruppe $V_4$ bekannt. Die Bezeichnung \emph{Vierergruppe} findet sich bei Felix Klein in seinen \emph{Vorlesungen über das Ikosaeder} von 1884 \cite{Klein1884}. In unserem Kontext erscheint $V_4$ nicht durch Geometriesymmetrien eines Rechtecks, sondern als Einheitengruppe
\[
(\Z/12\Z)^\times=\{1,5,7,11\}\simeq C_2\times C_2.
\]
Diese Identifikation ist elementar, aber für die Vierstellenstruktur eines Primzahlvierlings besonders passend.

\subsection{Hardy--Littlewood und lokale Zulässigkeit}

Hardy und Littlewood formulierten in den 1920er Jahren eine weitreichende Heuristik für Primzahlkonstellationen \cite{HardyLittlewood1923}. Für eine endliche Verschiebungsmenge $\mathcal H$ mit $k$ Elementen sei $\nu_{\mathcal H}(\ell)$ die Zahl der von $\mathcal H$ belegten Restklassen modulo einer Primzahl $\ell$. Die zugehörige singuläre Reihe lautet formal
\[
\mathfrak S(\mathcal H)
 =\prod_{\ell\ \mathrm{prim}}
 \frac{1-\nu_{\mathcal H}(\ell)/\ell}{(1-1/\ell)^k}.
\]
Für $\mathcal H=\{0,2,6,8\}$ gilt $k=4$. Lokale Zulässigkeit bedeutet $\nu_{\mathcal H}(\ell)<\ell$ für jede Primzahl $\ell$.

Unsere Zentralzahlformulierung ist eine symmetrische Version derselben lokalen Information. Die vier Primstellen lauten
\[
M-4,\quad M-2,\quad M+2,\quad M+4.
\]
Für jede Primzahl $\ell$ definieren wir daher die zulässigen Mittelpunkte
\[
\mathcal A_\ell
 =\bigl\{m\in\Z/\ell\Z:\;m\pm2\not\equiv0,\;m\pm4\not\equiv0\pmod\ell\bigr\}.
\]
Für $\ell=3,5,7$ erhält man unmittelbar
\[
\mathcal A_3=\{0\},\qquad
\mathcal A_5=\{0\},\qquad
\mathcal A_7=\{-1,0,+1\}.
\]
Gerade diese drei lokalen Aussagen bilden später den mod-$210$-Träger.

\section{Primzahlvierlinge und ihre Zentralzahl}

\begin{definition}[Primzahlvierling und Zentrum]
Für $p\in\N$ heiße
\[
Q(p)=(p,p+2,p+6,p+8)
\]
ein \emph{Primzahlvierling}, wenn alle vier Einträge prim sind. Seine Zentralzahl ist
\[
M:=p+4,
\]
sodass
\[
Q(p)=(M-4,M-2,M+2,M+4).
\]
\end{definition}

Der Vierling $Q(5)=(5,7,11,13)$ ist für die mod-$30$-Struktur ein kleiner Ausnahmefall. Im Folgenden heißt ein Primzahlvierling \emph{regulär}, wenn $p>5$.

\begin{lemma}[Lokale Bedingungen bei $2,3,5$]\label{lem:235}
Sei $Q(p)$ ein regulärer Primzahlvierling mit Zentralzahl $M$. Dann gilt
\[
M\equiv1\pmod2,\qquad M\equiv0\pmod3,\qquad M\equiv0\pmod5.
\]
Insbesondere
\[
\boxed{M\equiv15\pmod{30}.}
\]
\end{lemma}

\begin{proof}
Da $p>5$ prim ist, ist $p$ ungerade; damit ist auch $M=p+4$ ungerade. Modulo $3$ dürfen weder $M-4\equiv M-1$ noch $M-2\equiv M+1$ durch $3$ teilbar sein. Die Restklassen $M\equiv\pm1\pmod3$ sind daher ausgeschlossen, also $M\equiv0\pmod3$.

Modulo $5$ sind die vier verbotenen Werte
\[
M\equiv\pm2,\pm4\pmod5
\]
genau die vier von Null verschiedenen Restklassen. Daher bleibt nur $M\equiv0\pmod5$. Die drei Kongruenzen bestimmen durch den chinesischen Restsatz eindeutig $M\equiv15\pmod{30}$.
\end{proof}

\begin{lemma}[Siebung modulo $7$]\label{lem:7}
Für einen regulären Primzahlvierling gilt
\[
\boxed{M\pmod7\in\{-1,0,+1\}.}
\]
\end{lemma}

\begin{proof}
Eine Teilbarkeit einer der vier Zahlen $M\pm2,M\pm4$ durch $7$ ist ausgeschlossen. Somit sind die vier Mittelpunktklassen
\[
M\equiv\pm2,\pm4\pmod7
\]
verboten. Von den sieben Klassen bleiben genau $-1,0,+1$.
\end{proof}

\begin{satz}[Dreizellenstruktur modulo $210$]\label{thm:210}
Für jeden regulären Primzahlvierling gilt genau eine der drei Kongruenzen
\[
\boxed{M\equiv15,\ 105,\ 195\pmod{210}.}
\]
Äquivalent gilt für den Anfangsprimwert $p=M-4$
\[
\boxed{p\equiv11,\ 101,\ 191\pmod{210}.}
\]
\end{satz}

\begin{proof}
Nach Lemma~\ref{lem:235} ist $M\equiv15\pmod{30}$, nach Lemma~\ref{lem:7} ist $M\equiv1,0,-1\pmod7$. Da $\gcd(30,7)=1$, liefert der chinesische Restsatz genau drei Klassen modulo $210$. Direkte Einsetzung ergibt $15,105,195$.
\end{proof}

\section{Die EABC-Faser als Kleinsche Vierergruppe}

\begin{definition}[EABC-Sektoren]
Wir bezeichnen die vier Einheitenklassen modulo $12$ durch
\[
E\equiv1,\qquad A\equiv5,\qquad B\equiv7,\qquad C\equiv11\pmod{12}.
\]
\end{definition}

Multiplikativ bilden diese Klassen die Gruppe
\[
(\Z/12\Z)^\times=\{1,5,7,11\}\simeq V_4.
\]
Die Zerlegung folgt auch direkt aus
\[
\Z/12\Z\simeq \Z/4\Z\times\Z/3\Z
\]
und damit
\[
(\Z/12\Z)^\times\simeq(\Z/4\Z)^\times\times(\Z/3\Z)^\times\simeq C_2\times C_2.
\]

Wir definieren auf zu $6$ teilerfremden ganzen Zahlen
\[
\chi_4(n)=
\begin{cases}
+1,&n\equiv1\pmod4,\\
-1,&n\equiv3\pmod4,
\end{cases}
\qquad
\chi_3(n)=
\begin{cases}
+1,&n\equiv1\pmod3,\\
-1,&n\equiv2\pmod3.
\end{cases}
\]
Dann ergibt sich die Charaktertabelle
\begin{equation}\label{eq:char-table}
\begin{array}{c|cc}
 &\chi_4&\chi_3\\ \hline
E=1 &(+1)&(+1)\\
A=5 &(+1)&(-1)\\
B=7 &(-1)&(+1)\\
C=11&(-1)&(-1)
\end{array}
\end{equation}
und damit eine explizite Identifikation von $V_4$ mit $\{\pm1\}^2$.

\subsection{Das Orientierungsbit}

Für die vier Offsets
\[
d\in\{-4,-2,+2,+4\}
\]
setzen wir
\[
\sigma=(+1,-1,-1,+1),\qquad
\tau=(-1,+1,-1,+1).
\]
Ferner definieren wir für ungerades $M$
\[
\eps(M)=
\begin{cases}
+1,&M\equiv1\pmod4,\\
-1,&M\equiv3\pmod4.
\end{cases}
\]

\begin{satz}[Charakterfaktorisierung]\label{thm:char-factor}
Sei $M\equiv0\pmod3$ ungerade. Dann gilt für alle $d\in\{-4,-2,+2,+4\}$
\[
\boxed{\chi_4(M+d)=\eps(M)\,\sigma(d),}
\]
\[
\boxed{\chi_3(M+d)=\tau(d),}
\]
und daher mit $\chi_{12}=\chi_4\chi_3$
\[
\boxed{\chi_{12}(M+d)=\eps(M)\,\sigma(d)\tau(d).}
\]
\end{satz}

\begin{proof}
Modulo $3$ ist $M\equiv0$, daher hängen die Werte von $\chi_3(M+d)$ ausschließlich von $d\pmod3$ ab; die vier Werte sind $(-,+,-,+)$. Modulo $4$ lassen die Verschiebungen $\pm4$ die Klasse von $M$ unverändert, während $\pm2$ sie zwischen $1$ und $3$ vertauschen. Daraus folgt das Muster $\eps(M)(+,-,-,+)$.
\end{proof}

\begin{definition}[Orientierungen]
Ein Primzahlvierling heißt \emph{$\ABCE$-orientiert}, wenn seine vier Primstellen modulo $12$ die Folge
\[
(5,7,11,1)
\]
bilden. Er heißt \emph{$\CEAB$-orientiert}, wenn die Folge
\[
(11,1,5,7)
\]
auftritt.
\end{definition}

\begin{korollar}[Orientierung aus der Mitte]\label{cor:orientation}
Für jeden regulären Primzahlvierling gilt
\[
\boxed{M\equiv1\pmod4\iff \ABCE,}
\qquad
\boxed{M\equiv3\pmod4\iff \CEAB.}
\]
\end{korollar}

\begin{proof}
Setzt man $\eps(M)=+1$ beziehungsweise $-1$ in Satz~\ref{thm:char-factor} und liest die vier Paare $(\chi_4,\chi_3)$ über \eqref{eq:char-table} zurück, entstehen genau die beiden angegebenen Folgen.
\end{proof}

\section{Die Kreis-Sehnen-Kodierung}

\subsection{Das reguläre Kreismodell}

Sei $M\ge1$ und
\[
N:=2M+1.
\]
Wir betrachten ein reguläres $N$-Eck auf dem Einheitskreis. Die Ecken werden im Uhrzeigersinn mit den zyklischen Indizes
\[
0,1,\ldots,2M
\]
nummeriert und durch die zentrierte Beschriftungsfunktion
\[
\lambda_M(r)=
\begin{cases}
r,&0\le r\le M,\\
r-N,&M<r\le2M
\end{cases}
\]
mit den ganzen Zahlen
\[
0,1,\ldots,M,-M,\ldots,-1
\]
versehen. Der Punkt $0$ liegt am oberen Scheitel; positive Beschriftungen liegen auf der rechten, negative auf der linken Kreishälfte.

Für zwei verschiedene Ecken $r,s$ sei $[r,s]$ die zugehörige Sehne. Bei einem regulären $N$-Eck ist die Sehnenrichtung modulo $\pi$ durch die Summe $r+s\pmod N$ bestimmt. Daher definieren wir die Parallelklasse
\[
\mathcal C_c
=\bigl\{\{r,s\}:r\ne s,\ r+s\equiv c\pmod N\bigr\}.
\]
Da $N$ ungerade ist, besitzt die Kongruenz $2r\equiv c\pmod N$ genau eine Lösung; dieser degenerierte Fall $r=s$ wird entfernt. Die übrigen $2M$ Ecken paaren sich zu genau $M$ Sehnen.

Für eine Sehne sei $r_R$ der geometrisch weiter rechts liegende und $r_L$ der weiter links liegende Endpunkt. Wir setzen
\[
\delta_M([r,s])=\lambda_M(r_R)-\lambda_M(r_L)
\]
und definieren die gerichtete Parallelklassensumme
\[
D_M(c):=\sum_{\gamma\in\mathcal C_c}\delta_M(\gamma).
\]

\begin{figure}[t]
\centering
\begin{tikzpicture}[scale=1.05]
  \def\M{5}
  \def\N{11}
  \draw[thick] (0,0) circle (2.5);
  \foreach \i in {0,...,10} {
    \pgfmathsetmacro{\ang}{90-360*\i/11}
    \coordinate (v\i) at (\ang:2.5);
    \fill (v\i) circle (1.4pt);
  }
  \node[above] at (v0) {$0$};
  \node[right] at (v1) {$1$};
  \node[right] at (v2) {$2$};
  \node[right] at (v3) {$3$};
  \node[right] at (v4) {$4$};
  \node[below right] at (v5) {$5$};
  \node[below left] at (v6) {$-5$};
  \node[left] at (v7) {$-4$};
  \node[left] at (v8) {$-3$};
  \node[left] at (v9) {$-2$};
  \node[left] at (v10) {$-1$};
  % c=2 chords, excluding fixed vertex 1
  \draw[blue!60!black, line width=.9pt] (v0)--(v2);
  \draw[blue!60!black, line width=.9pt] (v3)--(v10);
  \draw[blue!60!black, line width=.9pt] (v4)--(v9);
  \draw[blue!60!black, line width=.9pt] (v5)--(v8);
  \draw[blue!60!black, line width=.9pt] (v6)--(v7);
  \node[blue!60!black] at (3.7,1.8) {$c=2$};
\end{tikzpicture}
\caption{Schematisches Kreismodell für $M=5$ und $N=11$. Die blau eingezeichneten Sehnen bilden die Parallelklasse $c=2$, denn für ihre Endpunktindizes gilt jeweils $r+s\equiv2\pmod{11}$. Der eindeutige Fixpunkt der Gleichung $2r\equiv2$ ist keine Sehne.}
\label{fig:circle}
\end{figure}

\subsection{Geschlossene Formel}

\begin{satz}[Kreis-Sehnen-Formel]\label{thm:Dformula}
Für $1\le c\le M$ gilt
\[
\boxed{
D_M(c)=M^2-(c-1)M-\Big\lfloor\frac c2\Big\rfloor.
}
\]
\end{satz}

\begin{proof}
Wir unterscheiden gerade und ungerade Richtungen.

Sei zunächst $c=2h$. Die $M$ Sehnen zerfallen in drei Gruppen: Sehnen vollständig auf der rechten Kreishälfte, Sehnen zwischen rechter und linker Hälfte und Sehnen vollständig auf der linken Hälfte. Die jeweiligen Summen der gerichteten Differenzen sind
\[
h(h+1),\qquad (M-2h)(M+1),\qquad -h^2.
\]
Die erste Summe ist die Summe der geraden Zahlen $2,4,\ldots,2h$; die dritte ist das Negative der Summe der ersten $h$ ungeraden Zahlen. Für die kreuzenden Sehnen ergibt die explizite Paarung der Indizes die mittlere Summe $(M-2h)(M+1)$. Insgesamt
\begin{align*}
D_M(2h)
&=h(h+1)+(M-2h)(M+1)-h^2\\
&=M^2-(2h-1)M-h.
\end{align*}
Dies ist die behauptete Formel, da $\lfloor(2h)/2\rfloor=h$.

Sei nun $c=2h+1$. Die drei analogen Teilsummen lauten
\[
(h+1)^2,\qquad (M-2h-1)(M+1),\qquad -h(h+1).
\]
Daraus folgt
\begin{align*}
D_M(2h+1)
&=(h+1)^2+(M-2h-1)(M+1)-h(h+1)\\
&=M^2-2hM-h,
\end{align*}
und wiederum stimmt dies mit der Formel überein, weil $\lfloor(2h+1)/2\rfloor=h$.
\end{proof}

\begin{korollar}[Nachbarrichtungen]\label{cor:adjacent}
Für $c\ge1$ mit $c+1\le M$ gilt
\[
\boxed{
D_M(c)-D_M(c+1)=
\begin{cases}
M,&c\text{ gerade},\\
M+1,&c\text{ ungerade}.
\end{cases}}
\]
\end{korollar}

\begin{proof}
Subtraktion der beiden Ausdrücke aus Satz~\ref{thm:Dformula} liefert
\[
M+\left\lfloor\frac{c+1}{2}\right\rfloor-\left\lfloor\frac c2\right\rfloor,
\]
was je nach Parität von $c$ gleich $M$ oder $M+1$ ist.
\end{proof}

\begin{satz}[Fundamentale Decoderidentität]\label{thm:decoder}
Für $M\ge3$ gilt
\[
\boxed{M=D_M(2)-D_M(3).}
\]
Insbesondere gilt für jeden Modul $q\ge1$
\[
\boxed{M\equiv D_M(2)-D_M(3)\pmod q.}
\]
\end{satz}

\begin{proof}
Dies ist der gerade Fall $c=2$ von Korollar~\ref{cor:adjacent}. Explizit
\[
D_M(2)=M^2-M-1,\qquad D_M(3)=M^2-2M-1.
\]
Die Differenz ist $M$.
\end{proof}

\begin{bemerkung}[Informationsgehalt]
Satz~\ref{thm:decoder} ist zugleich eine positive und eine negative Aussage. Positiv ist die Kreis-Sehnen-Kodierung verlustfrei: Aus zwei Richtungen wird $M$ exakt rekonstruiert. Negativ folgt daraus, dass der vollständige Vektor $\bigl(D_M(c)\bigr)_c$ keine von $M$ unabhängige Primzahlinformation enthalten kann. Jede weitergehende statistische Aussage muss daher aus der Verteilung der zugrunde liegenden Primzahlvierlinge auf Restklassen stammen, nicht aus einer Informationsvermehrung durch die Geometrie.
\end{bemerkung}

\section{Die \texorpdfstring{$3\times2$}{3x2}-Struktur modulo 420}

Die Dreizelle aus Satz~\ref{thm:210} wird durch das Orientierungsbit $M\pmod4$ verfeinert. Da $210\equiv2\pmod4$, besitzt jede Klasse modulo $210$ genau zwei Lifts modulo $420$ mit entgegengesetztem mod-$4$-Wert.

\begin{satz}[Sechs Zellen modulo $420$]\label{thm:420}
Für reguläre Primzahlvierlinge treten genau die folgenden sechs zulässigen Mittelpunktklassen auf:
\[
\boxed{15,225,105,315,195,405\pmod{420}.}
\]
Die Zuordnung zu Siebzelle und EABC-Orientierung ist
\[
\begin{array}{c|cc}
M\pmod7 & \ABCE\;(M\equiv1\!\pmod4) & \CEAB\;(M\equiv3\!\pmod4)\\ \hline
+1&225&15\\
0&105&315\\
-1&405&195
\end{array}
\]
jeweils modulo $420$.
\end{satz}

\begin{proof}
Die drei Klassen modulo $210$ sind $15,105,195$. Ihre beiden Lifts modulo $420$ sind jeweils $r$ und $r+210$. Die mod-$4$-Werte und Korollar~\ref{cor:orientation} geben die Tabelle.
\end{proof}

Damit wird die lokale Struktur in der Form
\[
\boxed{
\{+1,0,-1\}_{M\bmod7}\times\{+1,-1\}_{\eps(M)}
}
\]
sichtbar. Es handelt sich um eine $3\times2$-Zerlegung: drei Siebzellen und ein binäres Orientierungsbit.

\section{Dyadische Lifts}

\begin{satz}[Exakter Lift]\label{thm:lift}
Für ganze $Q$ gilt
\[
\boxed{
D_{M+Q}(c)-D_M(c)=Q\bigl(2M+Q-c+1\bigr).
}
\]
\end{satz}

\begin{proof}
Einsetzen von $M+Q$ in Satz~\ref{thm:Dformula} und Subtraktion; der Boden-Term hängt nur von $c$ ab und kürzt sich weg.
\end{proof}

Setze nun
\[
Q_k:=210\cdot2^k.
\]
Für eine reguläre Mittelpunktzahl ist $M$ ungerade und $Q_k$ gerade. Daher ist
\[
2M+Q_k-c+1
\]
gerade genau dann, wenn $c$ ungerade ist.

\begin{korollar}[Dyadischer Richtungsdetektor]\label{cor:dyadic}
Modulo $2Q_k$ gilt
\[
\boxed{
D_{M+Q_k}(c)-D_M(c)\equiv
\begin{cases}
0,&c\text{ ungerade},\\
Q_k,&c\text{ gerade}
\end{cases}
\pmod{2Q_k}.}
\]
\end{korollar}

Damit sind ungerade Richtungen für das bei $Q_k\to2Q_k$ neu hinzukommende binäre Lift-Bit blind, während gerade Richtungen dieses Bit unterscheiden. Die Zahl der durch die lokalen Bedingungen bei $2,3,5,7$ zugelassenen Mittelpunktklassen wächst entlang
\[
210,420,840,1680,\ldots
\]
formal wie
\[
3,6,12,24,\ldots=3\cdot2^k.
\]
Dies ist eine Aussage über zulässige Restklassen, nicht darüber, ob oder wie häufig jede Klasse tatsächlich Primzahlvierlinge enthält.

\section{Hadamard- und Dirichlet-Projektionen}

Die vier Positionsvektoren
\[
\begin{aligned}
\mathbf 1&=(1,1,1,1), & \sigma&=(1,-1,-1,1),\\
\tau&=(-1,1,-1,1), & \sigma\tau&=(-1,-1,1,1)
\end{aligned}
\]
bilden eine orthogonale Hadamard-Basis von $\mathbb R^4$. Mit
\[
\langle f,g\rangle=\frac14\sum_{d\in\{-4,-2,2,4\}}f(d)g(d)
\]
folgt aus Satz~\ref{thm:char-factor} sofort:

\begin{proposition}[Charakterprojektionen]\label{prop:proj}
Für einen regulären Primzahlvierling gilt
\[
\boxed{\langle\sigma,\chi_4\rangle=\eps(M),}
\qquad
\boxed{\langle\tau,\chi_3\rangle=1,}
\qquad
\boxed{\langle\sigma\tau,\chi_{12}\rangle=\eps(M),}
\]
während die entsprechenden Kreuzprojektionen verschwinden.
\end{proposition}

Der $\chi_3$-Kanal ist somit ein starrer Positionskanal. Die Kanäle $\chi_4$ und $\chi_{12}$ tragen das Orientierungsbit.

Für eine der drei Anfangsprimzahlzellen $r\in\{11,101,191\}\pmod{210}$ seien
\[
N_r=N_{r,\ABCE}+N_{r,\CEAB},\qquad
\Delta_r=N_{r,\ABCE}-N_{r,\CEAB}.
\]
Dann ist $N_r$ die Population der Zelle und $\Delta_r$ ihr Orientierungsbias. Summiert man Proposition~\ref{prop:proj} über alle Vierlinge einer Zelle, erhält man
\[
\sum_{Q\in r}\sum_d\tau(d)\chi_3(M+d)=4N_r,
\]
\[
\sum_{Q\in r}\sum_d\sigma(d)\chi_4(M+d)=4\Delta_r,
\]
und analog für $\chi_{12}$.

Diese Identitäten sind exakt. Ob $\Delta_r(X)$ bei wachsendem $X$ systematische Vorzeichen- oder Größenmuster zeigt, ist hingegen eine statistische Frage und nicht durch die algebraische Zerlegung entschieden.

\section{Numerische Illustration}

Zur Illustration wurde für $p\le X$ mit einem elementaren Eratosthenes-Sieb geprüft, wann
\[
p,\ p+2,\ p+6,\ p+8
\]
prim sind. Die folgenden Zahlen sind reproduzierbare endliche Zählwerte; sie werden nicht als asymptotische Aussage interpretiert.

\begin{table}[h]
\centering
\caption{Primzahlvierlinge bis $p\le X$ und Zellbias $\Delta_r=N_{r,\ABCE}-N_{r,\CEAB}$. Der Ausnahmevierling $p=5$ ist in der Gesamtzahl und der Gesamt-Orientierung enthalten, nicht jedoch in den drei regulären mod-$210$-Zellen.}
\label{tab:counts}
\begin{tabular}{r r r r c r}
\toprule
$X$ & $N(X)$ & $N_{\ABCE}$ & $N_{\CEAB}$ & $(\Delta_{11},\Delta_{101},\Delta_{191})$ & $\sum\Delta_r$\\
\midrule
$10^6$ & 166 & 84 & 82 & $(-1,2,0)$ & 1\\
$2\cdot10^6$ & 295 & 152 & 143 & $(1,7,0)$ & 8\\
$10^7$ & 899 & 450 & 449 & $(4,8,-12)$ & 0\\
\bottomrule
\end{tabular}
\end{table}

Für $X=10^7$ ergibt sich detailliert
\[
\begin{array}{c|rrrr}
p\pmod{210}&N_r&N_{r,\ABCE}&N_{r,\CEAB}&\Delta_r\\ \hline
11&318&161&157&4\\
101&288&148&140&8\\
191&292&140&152&-12
\end{array}
\]
und damit $\Delta_{11}+\Delta_{101}+\Delta_{191}=0$. Der einzelne Ausnahmevierling $p=5$ ist $\ABCE$-orientiert; deshalb lautet die Gesamtbilanz bis $10^7$ dennoch $450:449$.

Die exakte Aufhebung bei diesem speziellen Cutoff ist ein endliches numerisches Phänomen. Aus ihr folgt weder ein Symmetriesatz für alle $X$ noch eine asymptotische Gleichverteilung.

\section{Lean-4-Formalisierungsschema}

Die elementaren Teile der Arbeit eignen sich gut für eine Formalisierung in Lean 4 \cite{deMouraUllrich2021}. Die sinnvollste Beweisarchitektur besteht aus kleinen Lemmata statt eines monolithischen Hauptsatzes:

\begin{enumerate}
\item Definition von \texttt{PrimeQuadruplet p} und \texttt{center p := p + 4};
\item Herleitung der lokalen Kongruenzen modulo $2,3,5,7$ aus der Primzahligkeit;
\item Zusammenbau der drei Klassen modulo $210$ durch lineare Ganzzahlarithmetik;
\item Definition der Prädikate \texttt{ABCE} und \texttt{CEAB} und Nachweis ihrer Äquivalenz zu $M\equiv1$ bzw. $3\pmod4$;
\item Definition der algebraischen Funktion
\[
D(M,c)=M^2-(c-1)M-\lfloor c/2\rfloor
\]
und maschineller Beweis von $D(M,2)-D(M,3)=M$;
\item in einer zweiten Formalisierungsstufe Definition des tatsächlichen regulären $(2M+1)$-Ecks, der Parallelklassen als endliche Mengen von Sehnen und Beweis, dass deren endliche Summe mit der geschlossenen Formel übereinstimmt.
\end{enumerate}

Ein kompakter Zielsatz kann mathematisch als
\begin{align*}
\operatorname{PrimeQuadruplet}(p)\land p>5
\;\Longrightarrow\;&
M\pmod{210}\in\{15,105,195\}\\
&\land\ (\ABCE\lor\CEAB)\\
&\land\ M=D_M(2)-D_M(3)
\end{align*}
formuliert werden, mit $M=p+4$. Für die mod-$420$-Verfeinerung kommt die entsprechende Sechszellenaussage hinzu.

Die in dieser Arbeit angegebene Lean-Struktur ist ein Formalisierungsschema; eine vollständig geometrische, maschinengeprüfte Datei müsste zusätzlich die Sehnengeometrie selbst als endliche kombinatorische Struktur implementieren.

\section{Hauptsatz in zusammengefasster Form}

\begin{satz}[EABC-Kreis-Sehnen-Zerlegungs- und Decodersatz]\label{thm:main}
Sei
\[
Q(p)=(p,p+2,p+6,p+8)
\]
ein Primzahlvierling mit $p>5$ und Zentralzahl $M=p+4$. Dann gelten:

\begin{enumerate}[label=\textup{(\roman*)}]
\item \textbf{Lokaler Zentralträger:}
\[
M\equiv15\pmod{30}.
\]

\item \textbf{Dreizellen-Sieb:}
\[
M\pmod7\in\{-1,0,+1\},
\]
und daher
\[
M\equiv15,105,195\pmod{210}.
\]

\item \textbf{$V_4$-Faser und Orientierung:}
Die vier Primstellen liegen in
\[
(\Z/12\Z)^\times\simeq V_4,
\]
und
\[
M\equiv1\pmod4\iff\ABCE,
\qquad
M\equiv3\pmod4\iff\CEAB.
\]
Damit entstehen genau sechs zulässige Klassen modulo $420$.

\item \textbf{Kreis-Sehnen-Decoder:}
Für die gerichteten Parallelklassensummen des regulären $(2M+1)$-Ecks gilt
\[
D_M(c)=M^2-(c-1)M-\Big\lfloor\frac c2\Big\rfloor
\]
und insbesondere
\[
\boxed{M=D_M(2)-D_M(3).}
\]

\item \textbf{Dyadischer Lift:}
Für $Q_k=210\cdot2^k$ unterscheidet der Übergang $M\mapsto M+Q_k$ modulo $2Q_k$ das neu hinzukommende Bit genau in den geraden Sehnenrichtungen, während ungerade Richtungen dafür blind sind.
\end{enumerate}
\end{satz}

\begin{proof}
Die Aussagen (i) und (ii) folgen aus Lemma~\ref{lem:235}, Lemma~\ref{lem:7} und Satz~\ref{thm:210}; Aussage (iii) folgt aus Satz~\ref{thm:char-factor}, Korollar~\ref{cor:orientation} und Satz~\ref{thm:420}; Aussage (iv) folgt aus Satz~\ref{thm:Dformula} und Satz~\ref{thm:decoder}; Aussage (v) folgt aus Satz~\ref{thm:lift} und Korollar~\ref{cor:dyadic}.
\end{proof}

\section{Diskussion und Grenzen}

Die Konstruktion legt eine saubere Trennung mehrerer Ebenen nahe:
\[
\begin{aligned}
&\underbrace{M\equiv15\pmod{30}}_{\text{lokaler Zentralträger}}
\quad\times\quad
\underbrace{M\pmod7\in\{-1,0,+1\}}_{\text{Dreizellen-Sieb}}\\
&\hspace{5em}\times\quad
\underbrace{\eps(M)\in C_2}_{\text{EABC-Orientierung}}.
\end{aligned}
\]
Die vier äußeren Primstellen bilden zugleich eine $V_4$-Faser. Man sollte jedoch nicht sagen, $V_4$ sei von der $2,3$-Struktur unabhängig: Gerade
\[
(\Z/12\Z)^\times\simeq(\Z/4\Z)^\times\times(\Z/3\Z)^\times
\]
zeigt, dass die Vierergruppe selbst aus den lokalen Einheiteninformationen bei $2$ und $3$ entsteht. Die sinnvolle Trennung lautet daher: \emph{nichtinvertierbare Zentralbedingungen} versus \emph{invertierbare Charakterfaser der vier Primstellen}.

Ebenso ist der mod-$420$-Schritt keine neue Primzahltheorie. Er kombiniert die mod-$210$-Siebzelle mit dem bereits in $M$ enthaltenen mod-$4$-Bit. Der Kreisdecoder macht diese Information sichtbar und rekonstruiert sie exakt, erzeugt sie aber nicht.

Was offen und potenziell interessant bleibt, ist die \emph{Verteilung} der Primzahlvierlinge auf den sechs Zellen. Dafür bieten sich die Größen
\[
\Delta_r(X)=N_{r,\ABCE}(X)-N_{r,\CEAB}(X)
\]
als Charakter-bewertete Primzahlrennen an. Eine ernsthafte Untersuchung müsste Hardy--Littlewood-artige Erwartungswerte, lokale Faktoren, mögliche Chebyshev-Bias-Mechanismen und robuste Fehlerkontrollen einbeziehen. Die vorliegende Arbeit stellt hierfür vor allem die algebraische Koordinatisierung bereit.

\section{Ausblick}

Drei Fortsetzungen erscheinen natürlich:

\begin{enumerate}
\item \textbf{Analytische Verteilung:} Untersuchung der sechs mod-$420$-Zellen und ihrer Charakterdifferenzen mit Methoden für Primzahlrennen und $L$-Funktionen.
\item \textbf{Höhere lokale Primfaktoren:} Erweiterung der Mittelpunkt-Siebmenge $\mathcal A_\ell$ auf $\ell=11,13,\ldots$ und Vergleich mit der singulären Reihe des Primzahlvierlings.
\item \textbf{Vollständige Formalisierung:} Umsetzung der Kreisgeometrie als endliche Struktur in Lean 4 und maschineller Beweis der geschlossenen Sehnensummenformel.
\end{enumerate}

Die elementare Kernidentität
\[
M=D_M(2)-D_M(3)
\]
bleibt dabei der einfachste Ausdruck der Konstruktion: Zwei benachbarte gerichtete Sehnenrichtungen rekonstruieren die Zentralzahl und damit sämtliche lokalen Restklasseninformationen des betrachteten Primzahlvierlings.

\section*{Danksagung und Reproduzierbarkeit}

Die numerischen Tabellen beruhen auf einem deterministischen Primzahlsieb und können mit wenigen Zeilen Code reproduziert werden. Für eine arXiv-Einreichung empfiehlt es sich, den verwendeten Quellcode und eine Lean-Datei als ergänzendes Material beizulegen. Die vorliegende Fassung macht ausdrücklich keinen Anspruch auf einen Beweis der Unendlichkeit von Primzahlvierlingen oder auf eine Verbesserung der Hardy--Littlewood-$k$-Tupel-Heuristik.

\begin{thebibliography}{99}

\bibitem{Dirichlet1837}
P.~G.~L. Dirichlet,
\newblock Beweis des Satzes, dass jede unbegrenzte arithmetische Progression, deren erstes Glied und Differenz ganze Zahlen ohne gemeinschaftlichen Factor sind, unendlich viele Primzahlen enthält,
\newblock \emph{Abhandlungen der Königlichen Preußischen Akademie der Wissenschaften} (1837), 45--81.

\bibitem{Dickson1904}
L.~E. Dickson,
\newblock A new extension of Dirichlet's theorem on prime numbers,
\newblock \emph{Messenger of Mathematics} (2) \textbf{33} (1904), 155--161.

\bibitem{HardyLittlewood1923}
G.~H. Hardy and J.~E. Littlewood,
\newblock Some problems of ``Partitio Numerorum''; III: On the expression of a number as a sum of primes,
\newblock \emph{Acta Mathematica} \textbf{44} (1923), 1--70.

\bibitem{Klein1884}
F. Klein,
\newblock \emph{Vorlesungen über das Ikosaeder und die Auflösung der Gleichungen vom fünften Grade},
\newblock Teubner, Leipzig, 1884.

\bibitem{IwaniecKowalski2004}
H. Iwaniec and E. Kowalski,
\newblock \emph{Analytic Number Theory},
\newblock American Mathematical Society Colloquium Publications, Vol.~53, Providence, 2004.

\bibitem{HardyWright2008}
G.~H. Hardy and E.~M. Wright,
\newblock \emph{An Introduction to the Theory of Numbers}, 6th ed.,
\newblock revised by D.~R. Heath-Brown and J.~H. Silverman, Oxford University Press, 2008.

\bibitem{OEISA007530}
OEIS Foundation Inc.,
\newblock The On-Line Encyclopedia of Integer Sequences, A007530: Prime quadruples,
\newblock \url{https://oeis.org/A007530}, abgerufen am 7. Oktober 2026.

\bibitem{deMouraUllrich2021}
L. de Moura and S. Ullrich,
\newblock The Lean 4 theorem prover and programming language,
\newblock in: A. Platzer and G. Sutcliffe (eds.), \emph{Automated Deduction -- CADE 28}, LNCS 12699, Springer, 2021, 625--635.
\newblock DOI: \url{https://doi.org/10.1007/978-3-030-79876-5_37}.

\end{thebibliography}

\end{document}
